% !TEX root = ../main.tex
\section*{Data availability}

Source data, processed aggregate tables and figure-source files will be deposited in a versioned public archive before publication. The Tree-of-Life alignments are available from the studies in which they were published \cite{hug2016,williams2D}. HCV records are available from NCBI Nucleotide under accession versions NC\_004102.1, AJ238799.1, AB047639.1, D10988.1, D17763.1, D49374.1, Y11604.1, AF064490.1, Y12083.1 and EF108306.2; the accession bundle, alignments, complete 4,554-row window grid and figures are retained in the versioned HCV release. Raw simulation shards are retained on the Vienna Life Science Compute Cluster and are available from the corresponding author on reasonable request. The archival DOI and a reviewer-access link will be added before submission.

\section*{Code availability}

The development repository is hosted at \url{https://github.com/EnesSakalliUniWien/iqtree3}. The IQ-TREE integration, simulation and postprocessing scripts, plotting code and manuscript sources used for this study will be archived with the source data before submission. Exact source versions are recorded in the experiment and run manifests.

\begin{thebibliography}{99}

\bibitem{felsenstein1985confidence}
Felsenstein, J. Confidence limits on phylogenies: an approach using the bootstrap. \emph{Evolution} \textbf{39}, 783--791 (1985).

\bibitem{hoang2017}
Hoang, D. T., Chernomor, O., von Haeseler, A., Minh, B. Q. \& Vinh, L. S. UFBoot2: improving the ultrafast bootstrap approximation. \emph{Mol. Biol. Evol.} \textbf{35}, 518--522 (2018).

\bibitem{anisimova2006}
Anisimova, M. \& Gascuel, O. Approximate likelihood-ratio test for branches: a fast, accurate, and powerful alternative. \emph{Syst. Biol.} \textbf{55}, 539--552 (2006).

\bibitem{minh2020iqtree}
Minh, B. Q. et al. IQ-TREE 2: new models and efficient methods for phylogenetic inference in the genomic era. \emph{Mol. Biol. Evol.} \textbf{37}, 1530--1534 (2020).

\bibitem{philippe2011}
Philippe, H. \emph{et al.} Resolving difficult phylogenetic questions: why more sequences are not enough. \emph{PLoS Biol.} \textbf{9}, e1000602 (2011).

\bibitem{ho2004}
Ho, S. Y. W. \& Jermiin, L. S. Tracing the decay of the historical signal in biological sequence data. \emph{Syst. Biol.} \textbf{53}, 623--637 (2004).

\bibitem{xia2003index}
Xia, X., Xie, Z., Salemi, M., Chen, L. \& Wang, Y. An index of substitution saturation and its application. \emph{Mol. Phylogenet. Evol.} \textbf{26}, 1--7 (2003).

\bibitem{duchene2021}
Duch\^ene, D. A., Mather, N., Van Der Wal, C. \& Ho, S. Y. W. Excluding loci with substitution saturation improves inferences from phylogenomic data. \emph{Syst. Biol.} \textbf{71}, 676--689 (2021).

\bibitem{manuel2025satute}
Manuel, C., Sakalli, E., Schmidt, H. A., Viñas, C., von Haeseler, A. \& Elgert, C. When the past fades: detecting phylogenetic signal with SatuTe. \emph{Mol. Biol. Evol.} \textbf{42}, msaf090 (2025).

\bibitem{levin2017}
Levin, D. A. \& Peres, Y. \emph{Markov Chains and Mixing Times}. 2nd edn. American Mathematical Society (2017).

\bibitem{tavare1986}
Tavaré, S. Some probabilistic and statistical problems in the analysis of DNA sequences. \emph{Lect. Math. Life Sci.} \textbf{17}, 57--86 (1986).

\bibitem{hug2016}
Hug, L. A. \emph{et al.} A new view of the tree of life. \emph{Nat. Microbiol.} \textbf{1}, 16048 (2016).

\bibitem{williams2D}
Williams, T. A., Cox, C. J., Foster, P. G., Sz\H{o}ll\H{o}si, G. J. \& Embley, T. M. Phylogenomics provides robust support for a two-domains tree of life. \emph{Nat. Ecol. Evol.} \textbf{4}, 138--147 (2020).

\bibitem{woese1990}
Woese, C. R., Kandler, O. \& Wheelis, M. L. Towards a natural system of organisms: proposal for the domains Archaea, Bacteria, and Eucarya. \emph{Proc. Natl Acad. Sci. USA} \textbf{87}, 4576--4579 (1990).

\bibitem{baker2004}
Baker, G., Smith, J. J. \& Cowan, D. A. Review and re-analysis of domain-specific 16S primers. \emph{J. Microbiol. Methods} \textbf{55}, 541--555 (2004).

\bibitem{smith2014}
Smith, D. B. \emph{et al.} Expanded classification of hepatitis C virus into 7 genotypes and 67 subtypes: updated criteria and genotype assignment web resource. \emph{Hepatology} \textbf{59}, 318--327 (2014).

\bibitem{lemoine2018}
Lemoine, F. \emph{et al.} Renewing Felsenstein's phylogenetic bootstrap in the era of big data. \emph{Nature} \textbf{556}, 452--456 (2018).

\bibitem{rambaut1997}
Rambaut, A. \& Grass, N. C. Seq-Gen: an application for the Monte Carlo simulation of DNA sequence evolution along phylogenetic trees. \emph{Bioinformatics} \textbf{13}, 235--238 (1997).

\bibitem{yang2016}
Yang, B., Wang, Y. \& Qian, P.-Y. Sensitivity and correlation of hypervariable regions in 16S rRNA genes in phylogenetic analysis. \emph{BMC Bioinformatics} \textbf{17}, 135 (2016).

\end{thebibliography}

\section*{Acknowledgements}

This work builds on the SatuTe framework of Manuel and colleagues and on its integration with IQ-TREE. Computation used the Vienna Life Science Compute Cluster.

\section*{Funding}

Funding information will be completed before submission.

\section*{Author contributions}

B.S. conceived the extension, developed the statistical formulation, performed the analyses and wrote the manuscript.

\section*{Competing interests}

The author declares no competing interests.

\section*{Additional information}

Further information on research design is available in the Nature Portfolio Reporting Summary associated with this article.
